\subsection{Módulos cuasifinitos sobre los anillos locales}
\label{subsection:0.7.4-fr}

\begin{definition}[7.4.1]
\label{0.7.4.1-fr}
Dado un anillo local $A$, de ideal maximal $\mathfrak m$, se dice que un
$A$-módulo $M$ es cuasifinito (sobre $A$) si $M/\mathfrak mM$ es de rango finito
sobre el cuerpo residual $k=A/\mathfrak m$.
\end{definition}

Cuando $A$ es \emph{noetheriano}, el separado completado $\widehat M$ de $M$ para la
topología $\mathfrak m$-preádica es entonces un
$\widehat A$-\emph{módulo de tipo finito}; en efecto, como
$\mathfrak m/\mathfrak m^2$ es entonces un $A$-módulo de tipo finito, esto resulta
de (7.2.12) y de la hipótesis sobre $M/\mathfrak mM$.

En particular, si se supone además que $A$ es \emph{completo} y $M$
\emph{separado} para la topología $\mathfrak m$-preádica (es decir,
$\bigcap_n\mathfrak m^nM=0$), el propio $M$ es un
$A$-\emph{módulo de tipo finito}: en efecto, $\widehat M$ es entonces un
$A$-módulo de tipo finito, y como $M$ se identifica con un submódulo de
$\widehat M$, $M$ es también de tipo finito (y, por lo demás, idéntico a su
completado en virtud de (7.3.6)).

\begin{proposition}[7.4.2]
\label{0.7.4.2-fr}
Sean $A$, $B$ dos anillos locales, $\mathfrak m$, $\mathfrak n$ sus ideales
maximales, y supongamos que $B$ es noetheriano. Sean $\varphi:A\to B$ un
homomorfismo local, $M$ un $B$-módulo de tipo finito. Si $M$ es un $A$-módulo
cuasifinito, las topologías $\mathfrak m$-preádica y $\mathfrak n$-preádica
sobre $M$ son idénticas, por tanto separadas.
\end{proposition}

Observemos que por hipótesis $M/\mathfrak mM$ es de \emph{longitud finita}
como $A$-módulo, por tanto también \emph{a fortiori} como $B$-módulo. Se
concluye que $\mathfrak n$ es el \emph{único ideal primo} de $B$ que contiene
el anulador de $M/\mathfrak mM$: en efecto, se reduce inmediatamente, en virtud
de (1.7.4) y (1.7.2), al caso en que $M/\mathfrak mM$ es \emph{simple}, por tanto
necesariamente isomorfo a $B/\mathfrak n$, y nuestra afirmación es evidente
en este último caso. Por otra parte, como $M$ es un $B$-módulo de tipo finito,
los ideales primos que contienen el anulador de $M/\mathfrak mM$ son los
que contienen $\mathfrak mB+\mathfrak b$, designando por $\mathfrak b$
el anulador del $B$-módulo $M$ (1.7.5). Como $B$ es noetheriano, se concluye
([11], p.~127, cor.~4) que $\mathfrak mB+\mathfrak b$ es un ideal
\oldpage[0\textsubscript{I}]{69}
de definición de $B$, es decir, que existe $k>0$ tal que
$\mathfrak n^k\subset\mathfrak mB+\mathfrak b\subset\mathfrak n$; por consiguiente,
para todo $h>0$
\[
  \mathfrak n^{hk}M
  \subset(\mathfrak mB+\mathfrak b)^hM
  =\mathfrak m^hM
  \subset\mathfrak n^hM
\]
lo que prueba que las topologías $\mathfrak m$-preádica y
$\mathfrak n$-preádica sobre $M$ son las mismas; la segunda es, por lo demás,
separada en virtud de (7.3.5).

\begin{corollary}[7.4.3]
\label{0.7.4.3-fr}
Bajo las hipótesis de \textup{(7.4.2)}, si además $A$ es noetheriano y completo para
la topología $\mathfrak m$-preádica, $M$ es un $A$-módulo de tipo finito.
\end{corollary}

En efecto, $M$ es entonces separado para la topología $\mathfrak m$-preádica, y
nuestra afirmación resulta de la observación hecha en (7.4.1).

\begin{env}[7.4.4]
\label{0.7.4.4-fr}
El caso de aplicación de (7.4.2) que es más importante es aquel en que el propio $B$
es un $A$-\emph{módulo cuasifinito}, lo que equivale a decir que
$B/\mathfrak mB$ es una \emph{álgebra de rango finito sobre}
$k=A/\mathfrak m$; esta condición puede, por lo demás, descomponerse en la
conjunción de las dos siguientes, en virtud de lo que precede:
\begin{enumerate}
  \item[(i)] $\mathfrak mB$ \emph{es un ideal de definición de} $B$;
  \item[(ii)] $B/\mathfrak n$ \emph{es una extensión de rango finito del cuerpo}
  $A/\mathfrak m$.
\end{enumerate}

Cuando esto ocurre, todo $B$-módulo de tipo finito es evidentemente un
$A$-módulo cuasifinito.
\end{env}

\begin{corollary}[7.4.5]
\label{0.7.4.5-fr}
Bajo las hipótesis de \textup{(7.4.2)}, si $\mathfrak b$ es el anulador del
$B$-módulo $M$, $B/\mathfrak b$ es un $A$-módulo cuasifinito.
\end{corollary}

Supongamos $M\neq0$ (en caso contrario el corolario es evidente). Se puede considerar $M$
como un módulo sobre el anillo local noetheriano $B/\mathfrak b$; estando entonces su anulador
reducido a $0$, la demostración de (7.4.2) muestra que
$\mathfrak m(B/\mathfrak b)$ es un ideal de definición de
$B/\mathfrak b$. Por otra parte, $M/\mathfrak nM$ es un espacio vectorial de
rango finito sobre $A/\mathfrak m$, por ser un cociente de $M/\mathfrak mM$, que es
por hipótesis de rango finito sobre $A/\mathfrak m$; como $M\neq0$, se tiene
$M\neq\mathfrak nM$ en virtud del lema de Nakayama; como
$M/\mathfrak nM$ es un espacio vectorial $\neq0$ sobre $B/\mathfrak n$, el hecho
de que sea de rango finito sobre $A/\mathfrak m$ implica que $B/\mathfrak n$ es
también de rango finito sobre $A/\mathfrak m$; la conclusión resulta, pues, de (7.4.4)
aplicado al anillo $B/\mathfrak b$.

